\section{Introduction}
\label{sec:introduction}

\subsection{Motivation}
\label{sec:motivation}

Scour --- the removal of riverbed material from around bridge foundations by
flowing water --- is consistently identified as a leading cause of bridge
failure~\cite{wardhana2003analysis, arneson2012hec18}. For railway networks the
exposure is systemic: a single scour-critical crossing can close a line, and
probabilistic assessments of the British railway network attribute a
substantial share of its weather-related risk to bridge scour during
floods~\cite{lamb2019scour}. The hazard is also awkwardly timed. Scour develops
fastest during high flows, exactly when visual and diver inspections are most
dangerous and least reliable, and damage can refill with loose sediment before
an inspection takes place~\cite{prendergast2014review, arneson2012hec18}.
Instrumenting every vulnerable bridge permanently is rarely economical for a
large inventory, which motivates indirect approaches that reuse vehicles
already crossing the structure.

Drive-by monitoring extracts information about a bridge from the dynamic
response of a traversing vehicle, an idea introduced for bridge-frequency
identification by Yang et al.~\cite{yang2004extracting} and since developed
into a broad research field~\cite{malekjafarian2015review}, including dedicated
railway applications~\cite{braganca2023review, souza2023review}. The physical
basis for drive-by scour monitoring is established: foundation-stiffness loss
changes the soil--foundation and structural dynamics
\cite{prendergast2013investigation, prendergast2016determining}. Wavelet-based
analyses of simulated train-borne signals identified localized scour
signatures~\cite{fitzgerald2019drive}; vehicle pitch has shown specific scour
sensitivity~\cite{mcgeown2024pitch}; scour-induced stiffness loss has been
identified from field-measured drive-by data~\cite{tola2025scour}; and
machine-learning pipelines have classified scour severity under operational
variability~\cite{fernandes2024driveby, fernandes2025early,
fernandes2026early, fernandes2026canelas}. A vehicle-borne estimate of
foundation condition is also a natural input to decision-oriented monitoring
and digital-twin frameworks~\cite{kamariotis2023framework,
torzoni2024digital}, which benefit from a continuous component-level state
rather than a binary alarm.

\subsection{Related work: the drive-by scour and multi-damage line}
\label{sec:related}

This study continues the drive-by scour programme developed by Fernandes and
co-workers. Fernandes et al.~\cite{fernandes2024driveby} detected scour,
modeled as local pier-foundation stiffness loss, with an unsupervised deep
autoencoder operating on raw vertical accelerations. Comparing on-board
positions, they found the best detection at the bogies. Fernandes et
al.~\cite{fernandes2025early} then studied supervised multi-damage
classification of cracks, bearing-device damage, and scour using min--max
normalization, Piecewise Aggregate Approximation (PAA), and a CNN tuned by
Bayesian optimization. Fernandes et al.~\cite{fernandes2026early} classified
scour levels with a supervised CNN and showed the value of adding speed as an
input. Fernandes et al.~\cite{fernandes2026canelas} extended the approach to a
high-fidelity finite-element model of the Canelas railway bridge, calibrated to
field modal parameters, using a Bayesian-optimized CNN--LSTM.

Three wider strands shape the present design. First, drive-by learning studies
often hold the rail-irregularity profile fixed within a scenario while varying
speed, mass, and noise~\cite{locke2020using, fernandes2025early}; profile
handling must therefore be stated as an experimental choice. Second,
sprung-mass channels are filtered by the suspension, whereas unsprung or
axle-level responses act nearer the wheel--rail interface
\cite{molodova2011axle, malekjafarian2015review}, which motivates comparing
responses across suspension levels. Third, two-sensor residual concepts can
cancel shared profile content~\cite{keenahan2014use, obrien2015apparent}, so
pairs may contain information that neither member provides alone.

This literature establishes feasibility, but it does not give a controlled
answer to the linked design choices of representation, neural components, and
response channels. Published studies generally optimize one pipeline or a
small set of preselected positions. Independent searches can also mix method
quality with finite hyperparameter-search luck, while passage-level splitting
can leak the repeated observations of one bridge state across partitions.

\subsection{Gap and design questions}
\label{sec:gap}

The experiment reported here is prospectively specified and source-locked:
its scientific blocks, state catalogues, splits, training grid, selection
rules, and analysis roles were fixed in versioned, hash-identified source
before production data were generated. Throughout this paper, ``registered''
is shorthand for that repository-fixed specification; it denotes no external
registry deposit.

The settled Paper-1 campaign contains four blocks rather than a progressive
mechanism ladder. F40-S is a two-span, one-target scour benchmark; F40-M adds
nominal bearing fixity and local flexural-rigidity loss. L99-S and L99-M repeat
those scour-only and multi-damage roles on a four-span, three-target geometry.
All four blocks use one generated FRA-class-4 profile realization (registered
phase seed 20260728), shared by every state and passage. Speed, temperature,
and vehicle properties vary between passages.
Track-layer damage, wheel polygonization, and suspension damage are deferred by
configuration; their implementation remains in the one qualified source tree,
but they are not part of this experiment.

Within that scope, the study asks:

\begin{itemize}
  \item \textbf{RQ1.} On F40-S, how do RAW versus PAA input, Time2Vec-style
  spatial-coordinate encoding, LSTM recurrence, and fixed-width multi-rate
  pooling affect grouped-development scour error across a matched 16-cell HPO
  factorial?
  \item \textbf{RQ2.} Among the six learning-eligible modeled-DOF responses in
  the eight-response \texttt{physical8\_v1} payload,
  which unordered two-channel pair minimizes the prospectively specified
  aggregate screen score for the retained RAW/PAA pipelines, and what do the
  six non-selectable single-channel diagnostics show?
  \item \textbf{RQ3.} After independent selected-pair HPO in each block, how
  accurately do the retained pipelines estimate one or three support-specific
  stiffness-loss targets and localize the most affected support?
  \item \textbf{RQ4.} On the genuinely shared semantic states, how does the
  combined multi-damage physics/task block (nominal bearing fixity, local $EI$
  loss, and bearing heads) differ from its scour-only counterpart in F40 and
  L99?
  \item \textbf{RQ5.} How stable are the frozen configurations across a
  disjoint post-selection initialization set, and what descriptive change is
  observed when F40-S settings are transferred downstream without retuning?
\end{itemize}

Three claim boundaries apply throughout. The target is modeled vertical
support-stiffness loss, not physical scour-hole depth. The inputs are modeled
response channels, not transducer packages or placements. Specifically, the
two wheelset rows are idealized constrained-motion kinematic diagnostics, not
independent vehicle degrees of freedom or validated axle-box observations, and
they are excluded from every learning and selection input. Finally, each block
receives its own selected-pair HPO for the primary analysis, so cross-block
primary results do not measure frozen-hyperparameter transport.

\subsection{Contributions}
\label{sec:contributions}

The paper contributes the following experimental design and, once populated by
the authenticated campaign, its corresponding estimates:

\begin{enumerate}
  \item a four-block, state-grouped simulation design with explicit semantic
  pairing: a 30-state F40-S/F40-M controlled subset and complete
  475-state L99-S/L99-M pairing, with common random numbers only where those
  identities actually match;
  \item a balanced $2\times2\times2\times2$ architecture factorial covering
  RAW/PAA, spatial encoding, recurrence, and fixed-width multi-rate pooling,
  with five independent 100-trial HPO restarts per F40-S cell;
  \item a complete Option-C adjudication in which the five HPO winners per cell
  are compared on one predeclared three-fold grouped-development OOF partition
  before the outer test is opened;
  \item a prospectively listed \texttt{physical8\_v1} response screen over the
  six eligible, non-selectable singles and all 15 selectable pairs for four
  retained RAW/PAA pipeline slots, with paired initialization seeds and
  authenticated alias deduplication;
  \item an optional secondary contemporary-architecture tier admitted only
  after the F40-S screen, in which every challenger uses the same authenticated
  primary selected pair and no challenger performs its own pair selection;
  \item independent selected-pair HPO and disjoint post-freeze stability fits
  in every scientific block, separating primary block-local performance from a
  secondary frozen-F40-S transfer analysis; and
  \item explicit numerical and observation-model boundaries for the
  support-aligned meshes, inherited track topology and scaling, Rayleigh
  damping policy, and diagnostic-only idealized wheelset response proxies.
\end{enumerate}

Table~\ref{tab:delta} summarizes the relation to the closest prior work.
\pending{Add one sentence after Table~\ref{tab:delta} stating the headline
authenticated four-block outcome.}

\begin{table}[htbp]
\centering
\caption{This study in relation to the drive-by scour line of Fernandes et
al. Prior-work entries summarize the cited primary sources.}
\label{tab:delta}
\footnotesize
\setlength{\tabcolsep}{3pt}
\begin{tabular}{p{2.2cm}p{2.8cm}p{2.8cm}p{3.0cm}p{3.2cm}}
\toprule
Work & Task / mechanisms & Pipeline & Selection & Robustness / analysis \\
\midrule
Fernandes 2024~\cite{fernandes2024driveby} & unsupervised scour detection &
raw signal + deep autoencoder & fixed AE; on-board positions compared & KL
damage index, ROC \\
\addlinespace
Fernandes 2025~\cite{fernandes2025early} & classification of crack + bearing +
scour & min--max + PAA + CNN & Bayesian-optimized CNN; two fixed channels &
confusion matrices, boxplots, EOV sensitivity \\
\addlinespace
Fernandes 2026a~\cite{fernandes2026early} & classification of scour levels &
CNN (+ speed input) & Bayesian-optimized CNN; two fixed channels & confusion
matrices, boxplots \\
\addlinespace
Fernandes 2026b~\cite{fernandes2026canelas} & classification of scour depths
(field-calibrated model) & CNN--LSTM & Bayesian-optimized per channel
configuration & multi-run statistics \\
\addlinespace
\textbf{This work} & \textbf{continuous one-/three-support regression,
localization, and bearing co-estimation in four explicit blocks} &
\textbf{RAW/PAA $\times$ position $\times$ LSTM $\times$ fixed-width
multi-rate factorial} & \textbf{Option-C grouped OOF adjudication; complete
  6-single/15-pair screen; per-block selected-pair HPO} & \textbf{sealed outer
test; disjoint post-freeze stability; authenticated paired-state and secondary
transfer roles} \\
\bottomrule
\end{tabular}
\end{table}

\subsection{Paper organization}

Section~\ref{sec:framework} gives the experimental overview.
Section~\ref{sec:simulation} describes the TTBI model, the four state
catalogues, and active/deferred mechanisms. Section~\ref{sec:data_processing}
details preprocessing, architectures, HPO, adjudication, and response-channel
selection. Section~\ref{sec:results} provides results and discussion;
Section~\ref{sec:limitations} states claim boundaries and practical
feasibility; and Section~\ref{sec:conclusion} concludes.
